\chapter{Use Cases, Requirements and System Design}
\label{ch:usecases-design}

\Cref{ch:related-work} ended with the gap that the thesis addresses: when logic in
the loss pays, under a matched comparison, on several datasets and with its cost
measured. This chapter turns that gap into five use cases, requirements traceable
to the experiments, a reference architecture and the design of the knowledge base.
It does not fix how measurements are taken (\cref{ch:benchmarking-methodology}),
the code (\cref{ch:implementation}) or the numbers (\cref{ch:results}).
UC1, UC2 and UC3 are the tasks on which detectors are trained and scored on
labelled data. UC4 and UC5 are properties asked of the same detectors: that they
hold up on attacks they did not see, and that an alert comes with something a
human can read.

% ================================================================
% 7.1 UC1
% ================================================================
\section{Use case UC1 -- Batch multi-class traffic classification}
\label{sec:uc1}

In UC1 a security analyst holds a set of flows from a capture and wants each one
assigned to a class, Normal or an attack family of the dataset. The analysis is
offline: order does not matter and the time to produce labels is a cost, not a
constraint. This is the setting in which RQ1 is answered for every detector of the
comparison.

\subsection{Actors, inputs, outputs, success criteria}
\label{sec:uc1-actors}

The actors are the analyst, who reads the labels; the benchmark runner, which
trains each arm and scores it on held-out data; and the knowledge engineer, who
writes the axioms (\cref{sec:kb-design}). The arms are the detectors under
comparison: one network trained with cross-entropy, without and with class weights
(MLP-CE); the same network trained with LTN on the supervised axioms alone (LTN-sup)
or on the supervised and the domain axioms (LTN-full); Random Forest and XGBoost; and
the axioms alone, with nothing learned (rule-only). The input is a table of flow records,
reconstructed from captured traffic by a flow exporter \citep{claise2013rfc7011}.
The detector does not see packets. On the CIC datasets the rows come from
CICFlowMeter \citep{sharafaldin2018cicids2017}. NSL-KDD ships connection records
with content features computed when the dataset was built
\citep{tavallaee2009kddcup99}. Where a statement holds for both kinds of unit, this
chapter writes flow; on NSL-KDD the unit is a connection record. No payload is
inspected (\cref{ch:introduction}).
The classes of each dataset and the labels behind them are in \cref{tab:impl-classes}.
The output is a class for each flow with a score for every class.

The success criteria are measures, not thresholds. On quality: macro-F1, the
Matthews correlation coefficient, per-class precision and recall, and an alert
budget, the false alarms per 100\,000 benign flows. The budget is there because
attacks are rare in real traffic, so a small false positive rate can still bury
the analyst \citep{axelsson2000baserate}. On cost: time and memory to train and
run the detector. On usability: the effort of writing the knowledge and the
sensitivity of the result to it. A gain is claimed only against baselines of
matched capacity, under a protocol that avoids the pitfalls listed by
\citet{arp2022dosdonts}. UC1 runs on NSL-KDD (Normal, DoS, Probe, R2L, U2R) in
E2, and on the corrected CIC-IDS2017 \citep{engelen2021troubleshooting}, mapped to
the same families plus DDoS and Botnet, in E5.

% ================================================================
% 7.2 UC2
% ================================================================
\section{Use case UC2 -- DDoS detection in near real time}
\label{sec:uc2}

In UC2 an operator watches a link and wants to know, as flows complete, whether an
attack is under way. Detection is the first step of a mitigation workflow
\citep{dobbins2021rfc8903}. This is the use case for which the streaming pipeline
of \cref{sec:streaming-pipeline} exists, and it supplies the real-time half of
the answer to RQ3. The task is binary, Normal or DDoS. The input is a stream of
flows from a live capture or a replayed dataset. The output is a score for each
flow and an alert when the score reaches the operating point, a threshold fixed on
validation data.
The rule that fixes the operating point is given in \cref{sec:operating-point}.

Two assumptions are made explicit. First, flows arrive reassembled: a flow reaches
the detector when the exporter closes it, by end of connection, reset or timeout,
and the detector never sees a partial flow. Second, any feature computed over a
window of flows, such as the number of flows to one destination in the last
seconds, uses past flows only and has one definition shared by training and
streaming. A model trained on one definition and served with another would be
scored on a different task.

The binary task is a choice of UC2. CIC-DDoS2019 holds many attack types
\citep{cic2019ddos}, and the thesis merges them into one class, DDoS, because UC2
asks for detection of the attack and not for the name of the variant. The
consequence is declared now: the hypothesis on rare classes (H2, \cref{tab:hypotheses})
cannot be verified on this dataset, which has no rare class and no attack class
without an axiom. What it does
offer is a test day with an attack type, Portmap, absent from the training day,
the scenario of UC4. The experiments of UC2 are E6 in batch, and E14 and E15 in
streaming.

\subsection{Latency and throughput requirements, alert semantics}
\label{sec:uc2-latency}

We fix no absolute latency or throughput target: there is no deployment to take
one from, and an untested number would be a guess. We require instead that two
quantities be measured: the service time per flow and the throughput in flows per
second under a replayed load (E14), and the detection delay, from the start of an
attack to the first alert (E15). The model contributes its inference time, but we
expect the larger term in the delay to be the timeout of the flow exporter, which
belongs to the capture. The parts of the delay are reported separately, so that the share
of the model is visible.

An alert refers to one flow and carries the time the flow was closed, the time of
the decision, the score and, when a knowledge base is attached, the axioms that
hold on it (\cref{sec:uc5}). One correct alert tells the operator that an attack
has begun, while the alert budget counts every false alarm.

% ================================================================
% 7.3 UC3
% ================================================================
\section{Use case UC3 -- Rare and low-footprint attack detection (R2L / U2R)}
\label{sec:uc3}

R2L and U2R attacks leave little in flow-level features: the evidence is in the
data portion of a connection, which flow records do not carry
(\cref{sec:r2l-u2r-hard}). The classes are also small: NSL-KDD has 52 U2R and 995
R2L records in training, and 67 and 2\,885 in test (\cref{tab:impl-splits}).
Input and output are those of UC1. The success criteria are the recall and the
precision--recall area of each rare class, read along the chain of arms: MLP-CE
without class weights, MLP-CE with class weights, LTN-sup, LTN-full. Each step
changes one thing, the class balance, then the form of the loss, then the
knowledge, so that a gain can be placed. U2R is so small that the seeds are needed
to say anything.

We expect a negative result for the knowledge on NSL-KDD, and state it before the
experiment. The candidate axioms of \cref{ch:attacks-ddos} cover the volumetric
families, and with the features of the dataset no axiom for R2L or U2R was precise
enough to use (\cref{sec:kb-elicitation}). LTN-full receives no knowledge about
these classes and can differ from LTN-sup on them only through the effect of the
other axioms on the shared network. A variant with broad axioms for the two classes
tests whether weak knowledge helps or harms (\cref{sec:kb-format}). On CIC-IDS2017
a variant gives R2L, there FTP and SSH password guessing, one axiom; no knowledge
base has a precise axiom for U2R. UC3 is evaluated in E2, E3 and E5, and E8 varies the
imbalance.

% ================================================================
% 7.4 UC4
% ================================================================
\section{Use case UC4 -- Detection of previously unseen attack variants}
\label{sec:uc4}

In UC4 the analyst meets an attack that was not in the training data and wants an
alert anyway. The success criterion is the recall on the unseen attack at the
operating point fixed on validation data, with the quality on the whole test day.
Two scenarios serve it, both in E9: the Portmap attack of CIC-DDoS2019, present in
the test day and absent from the training day, and the transfer between the two CIC
datasets. In the second, attacks and flow-extractor version change together, so
the result measures both, and we say so where we read it.

Knowledge matters here in a specific way. An axiom that states the mechanism of an
attack family may hold on a variant never seen, while an axiom read off the
training data may not. This is why each axiom is tagged with its origin
(\cref{sec:kb-format}). We make no claim about the outcome.

% ================================================================
% 7.5 UC5
% ================================================================
\section{Use case UC5 -- Explainable alert generation for a human analyst}
\label{sec:uc5}

In UC5 the analyst triages alerts and needs, for each, a reason that can be read
and checked against the flow (\cref{sec:operator-needs}). The thesis takes the part
the symbolic component can supply: the class, the score, and the axioms whose
antecedent holds on the flow, with their degree and the values of the conditions.
An LTN-trained network is a plain network at inference
(\cref{sec:symbolic-placement}), so these axioms are a rule match computed beside
the network, not the reasons for its decision. Experiment E17 measures coverage,
the share of alerts with at least one axiom; faithfulness, the agreement between the
class the axiom implies and the class the network predicts; and conciseness, the
number of axioms that fire on an alert. Faithfulness exposes the gap, since an explanation that
agrees with the prediction only by chance is worth little. Alerts can be handed on
in an exchange format \citep{debar2007rfc4765,opentext2023cef,jordan2021stix}; the
thesis does not depend on one, and the alert sink takes a plain record.
% FIXME(open): the serialisation of the alert record (IDMEF, CEF, STIX or plain JSON) is fixed in Ch. 9, \cref{ch:implementation}.

% ================================================================
% 7.6 Requirements
% ================================================================
\section{Functional and non-functional requirements}
\label{sec:requirements}

A functional requirement says what the system does for the use cases of
\cref{sec:uc1,sec:uc2,sec:uc3,sec:uc4,sec:uc5}. A non-functional requirement says
how well it does so, or under what conditions. Each is written so that an
experiment can check it and a measure can report on it. Each has an identifier,
the experiment that checks it, and one measure on one of the three axes of the
thesis: detection quality (Q), computational cost (C) and usability (U), in batch
and in real time (\cref{ch:introduction}). The experiments are numbered E1 to
E17. This chapter names them and states what each is for; their protocol is the
subject of \cref{ch:benchmarking-methodology} and their outcomes of
\cref{ch:results}.

The requirements carry no numeric target: a target would need a deployment or a
result to calibrate it on, and we have neither. A requirement is verified when its
experiment has run and its measure is reported with the spread over seeds, against a
baseline of matched capacity where the outcome is a comparison. A requirement that
the knowledge arms fail is a result, and the thesis reports it (contribution C3).

\subsection{Requirement table with IDs, traceable to experiments and benchmarks}
\label{sec:requirement-table}

\Cref{tab:functional-requirements} lists the functional requirements and
\cref{tab:nonfunctional-requirements} the non-functional ones.

\begin{table}[t]
\centering
\footnotesize
\caption{Functional requirements. Axis: Q quality, C cost, U usability.}
\label{tab:functional-requirements}
\setlength{\tabcolsep}{4pt}
\begin{tabular}{@{}p{1cm}p{5.5cm}p{1.2cm}p{0.7cm}p{4.2cm}@{}}
\toprule
ID & Requirement & Exp. & Axis & Measure \\
\midrule
FR1 & Assign each flow to one class of the dataset and return a score per class (UC1)
  & E2, E5 & Q & Macro-F1 and MCC on the test part \\
FR2 & Decide Normal or DDoS for each flow and raise an alert at a fixed operating point (UC2)
  & E6 & Q & MCC; false positive rate at a fixed true positive rate; alert budget \\
FR3 & Report R2L and U2R separately from the other classes (UC3)
  & E2, E5, E8 & Q & Per-class recall and precision--recall area \\
FR4 & Detect an attack type absent from the training data (UC4)
  & E9 & Q & Recall on the unseen attack at the operating point fixed on validation data \\
FR5 & Attach to an alert the axioms whose antecedent holds on the flow (UC5)
  & E17 & U & Explanation coverage, faithfulness and conciseness \\
FR6 & Read the axioms from a versioned file, so that an axiom set can be swapped or reduced without a change of code (\cref{sec:kb-format})
  & E3 & U & Change in macro-F1 and per-class recall when an axiom is removed \\
FR7 & Report for any model, knowledge arm or not, the satisfaction of each axiom and the rate at which its prediction violates the rule
  & E2 & Q, U & Axiom-level satisfaction and domain-rule violation rate \\
FR8 & Allow an optional symbolic check on the predictions, outside the network (\cref{sec:symbolic-placement})
  & E12 & C & Added inference latency per flow; violation rate before and after \\
FR9 & Give the same decision to the same flow in batch and in streaming (UC2)
  & E14 & Q & Share of identical decisions on a replay \\
\bottomrule
\end{tabular}
\end{table}

\begin{table}[t]
\centering
\footnotesize
\caption{Non-functional requirements. Axis: Q quality, C cost, U usability.}
\label{tab:nonfunctional-requirements}
\setlength{\tabcolsep}{4pt}
\begin{tabular}{@{}p{1.1cm}p{5.4cm}p{1.2cm}p{0.7cm}p{4.2cm}@{}}
\toprule
ID & Requirement & Exp. & Axis & Measure \\
\midrule
NFR1 & Run the NSL-KDD pipeline under the corrected protocol, with validation kept apart from the official test file
  & E1 & Q & Accuracy and macro-F1 on validation and on test \\
NFR2 & Keep the quality when labels are scarce or the training set is small
  & E7 & Q & Macro-F1 against the fraction of labels and the size of the training set \\
NFR3 & Keep the quality on rare classes when the imbalance grows
  & E8 & Q & Per-class precision--recall area against the imbalance \\
NFR4 & Keep the quality under noise and simple evasion of the flow features
  & E10 & Q & Change in macro-F1 under the perturbation \\
NFR5 & Keep the benefit of the knowledge when the axioms are degraded
  & E11 & U & Macro-F1 of LTN-full minus LTN-sup against the degradation level \\
NFR6 & Make the effect of the logic configuration visible (exponent $p$, crisp or graded conditions, scale of the widths)
  & E4 & U & Spread of macro-F1 over the configurations \\
NFR7 & Report the cost of training for every arm
  & E13 & C & Time per epoch, epochs, total fit time, peak memory; time per epoch against the number of axioms \\
NFR8 & Keep the inference cost of an LTN-trained network equal to the same network trained with cross-entropy
  & E12 & C & Inference latency and throughput \\
NFR9 & Sustain the arrival rate of flows in streaming
  & E14 & C & Throughput in flows per second; latency quantiles \\
NFR10 & Detect a DDoS attack early
  & E15 & C & Time from the start of the attack to the first alert \\
NFR11 & Report the effort of writing and revising the axioms
  & E16 & U & Estimated hours per axiom, per base and per rule change; axioms and lines per base \\
NFR12 & Make a run reproducible from its configuration and its seed, with the version of the knowledge base recorded
  & E1, E2 & Q & Equality of the metrics of two runs of one seed on one machine \\
NFR13 & Summarise the usability of each arm on a scorecard whose anchors are written before the measurements
  & E16, E17 & U & Score per criterion and arm \\
\bottomrule
\end{tabular}
\end{table}

FR9 and NFR12 are our additions to the experiments of the plan, and cost almost
nothing to check: FR9 is a consistency check between the two pipelines, and NFR12 is
the condition under which any other row can be trusted. NFR13 is the rubric of the
usability axis announced in the introduction. E14 and E15 depend on the
streaming infrastructure; if it is not available to us, FR9, NFR9 and NFR10 are
specified and the measurements left open (\cref{sec:streaming-pipeline}). E1 and E2
use NSL-KDD, E5 CIC-IDS2017 and E6 CIC-DDoS2019; the datasets of the other
experiments are assigned in \cref{ch:benchmarking-methodology}. Rows on UC3 are run
only where a rare class exists, which excludes CIC-DDoS2019 (\cref{sec:uc2}). Every
row that compares arms covers the four neural arms, which share one network and one
batch size, and Random Forest, XGBoost and the rule-only arm. A one-dimensional
convolutional network enters the final runs and the experiments on saved models as a
further deep baseline, outside matched capacity.

% ================================================================
% 7.7 Reference architecture
% ================================================================
\section{Reference architecture}
\label{sec:reference-architecture}

The batch pipeline serves UC1, UC3 and UC4 and produces the results of the
comparison. The streaming pipeline serves UC2 and UC5 and produces the real-time
half of RQ3. They differ in where flows come from and when the model is called,
and share the model, the feature definitions and the knowledge base, so that a
result in one is a statement about the other (\cref{fig:pipelines}).

\begin{figure}[t]
\centering
\footnotesize
\setlength{\fboxsep}{3pt}
\newcommand{\pbox}[1]{\fbox{\parbox[c][0.9cm][c]{1.75cm}{\centering #1}}}
\textbf{Batch}\\[3pt]
\pbox{Capture or dataset} $\rightarrow$ \pbox{Flow extraction} $\rightarrow$
\pbox{Pre\-processing} $\rightarrow$ \pbox{Model} $\rightarrow$ \pbox{Report}\\[8pt]
\textbf{Streaming}\\[3pt]
\pbox{Capture or replay} $\rightarrow$ \pbox{Flow assembly} $\rightarrow$
\pbox{Feature vector} $\rightarrow$ \pbox{Broker}\\[3pt]
$\downarrow$\\[3pt]
\pbox{Inference service} $\rightarrow$ \fbox{\parbox[c][0.9cm][c]{1.75cm}{\centering
Symbolic check (optional)}} $\rightarrow$ \pbox{Alert sink} $\rightarrow$
\pbox{Dashboard}
\caption{The batch and the streaming pipeline. The symbolic check is present
only in the variants that ask for it (\cref{sec:symbolic-placement}).}
\label{fig:pipelines}
\end{figure}

\subsection{Batch pipeline: capture $\rightarrow$ flow extraction $\rightarrow$ preprocessing $\rightarrow$ model $\rightarrow$ report}
\label{sec:batch-pipeline}

A capture, or a dataset that already holds flows, is turned into a table of
flows by a flow extractor. Preprocessing encodes the categorical columns and
scales the numeric ones; it is fitted on the training part only and applied
unchanged to the validation and test parts. The model step trains one arm and
returns class scores for each part. The report step computes the measures of the
three axes and writes them, with the configuration, the seed and the version of
the knowledge base, to one record. A run is one configuration and one seed, and
every table of the thesis is built from such records, which is what makes NFR12
checkable. Which parts a run may touch is decided by the protocol of
\cref{ch:benchmarking-methodology}.

The axioms are written on the raw columns of the flow table, not on the scaled
features that the network receives. This keeps their thresholds in the units of
the traffic, where an analyst can read them, and it means the report step needs
both the raw table and the scaled features.

\subsection{Streaming pipeline: capture/\allowbreak replay $\rightarrow$ flow assembly $\rightarrow$ feature vector $\rightarrow$ broker $\rightarrow$ inference service $\rightarrow$ (symbolic check) $\rightarrow$ alert sink $\rightarrow$ dashboard}
\label{sec:streaming-pipeline}

The source is a live capture or the replay of a captured dataset, at its
recorded rate or a scaled one. Flow assembly closes the flows, so that they reach
the rest of the pipeline reassembled (\cref{sec:uc2}). The feature vector is built
with the code and definitions of training, window features included. A broker
decouples producer and consumer, absorbs bursts and lets the inference service run
at its own rate. The service loads the exported model and calls it on micro-batches.
An optional symbolic check follows; the alert sink collects the decisions that
reach the operating point, and a dashboard shows them.
% FIXME(open): the choice of the broker and the implementation of each stage are open in Ch. 9, \cref{sec:realtime-pipeline}, and belong to the co-supervisor.

Three points need a decision or a caution. The first is the extractor. The corrected
CIC-IDS2017 comes from a corrected CICFlowMeter and CIC-DDoS2019 from the original
one, and a model must be served with flows from the extractor that produced its
training features. For streaming the model should be the one trained on the corrected
CIC-IDS2017, with the corrected extractor live. Its classes are those of CIC-IDS2017,
and for UC2 an alert is raised on the DDoS class.
% FIXME(open): confirm which model and which extractor the streaming pipeline uses.
The second is that a symbolic check needs the raw flow record beside the feature
vector, so the message on the broker carries both. The third is time: the pipeline
stamps each flow when it is closed, reaches the broker, is decided and is alerted.
The differences give the parts of the detection delay (\cref{sec:uc2-latency}) and
are the data of E14 and E15.

The inference service can run in one process or on a distributed engine through a
distributed deep learning (DDL) framework. The thesis has not decided to use DDL: it
needs a cluster available to us and a framework that runs the exported model. If
either is missing, the reference pipeline is the single-process one, which is also
the setting of the scope in \cref{ch:introduction}; if DDL is used, that scope
sentence, which excludes distributed inference, must change.
% FIXME(open): DDL is undecided (cluster access and a framework that runs the exported model); if used, amend the scope paragraph of Ch. 1, Sec. 1.6.
The mandatory part is the single-process pipeline: design, assumptions and method
of measurement. Measurements on a cluster are an extension.

\subsection{Where the symbolic component lives in each variant}
\label{sec:symbolic-placement}

The symbolic component is the knowledge base and the evaluation of its axioms.
\Cref{tab:symbolic-placement} states where each arm uses it.

\begin{table}[t]
\centering
\small
\caption{Use of the knowledge base by each arm, in training and at inference.}
\label{tab:symbolic-placement}
\setlength{\tabcolsep}{4pt}
\begin{tabular}{@{}p{2.6cm}p{3.7cm}p{3.1cm}p{3cm}@{}}
\toprule
Arm & Training & Inference & Post-check (E12) \\
\midrule
MLP-CE, with or without weights & None & Network only & Possible \\
LTN-sup & Supervised axioms in the loss & Network only & Possible \\
LTN-full & Supervised and domain axioms in the loss & Network only & Possible \\
Rule-only & None; nothing is learned & The axioms are the model & Not applicable \\
Random Forest, XGBoost & None & Model only & Possible \\
\bottomrule
\end{tabular}
\end{table}

In the LTN arms the symbolic component acts through the loss, and only at
training: the class predicates are the outputs of the network, and the axioms
shape its parameters (\cref{ch:ltn}). At inference the trained network is a plain
multilayer perceptron with the same number of parameters and the same forward pass
as the baseline. In the rule-only arm the axioms are the whole model: conditions
are hard thresholds, the first axiom in the order of the file wins a conflict, and
a flow that no axiom covers goes to Normal.

The post-check is optional for every model that outputs a class. It evaluates the
knowledge base on the raw flow record after the model has decided, in one of two
modes. It annotates the alert with the axioms that fire, which is the source of the
explanations of UC5. Or it enforces the rules: when an axiom fires and the prediction
differs from its target, the prediction is replaced by the target, and a conflict
between axioms with different targets goes to the first in the order of the file, as
in the rule-only arm. The enforcing mode is the post-hoc symbolic filter that the
design keeps as an extension, a control for the neurosymbolic arms. It removes every
violation except those of an axiom that loses a conflict, and the interesting
quantity is what it costs in quality, because the axioms are not exact
(\cref{sec:kb-format}).

E12 has two parts. The first compares the inference latency and throughput of an
LTN-trained network with those of the same network trained with cross-entropy. They
are expected to be equal because the networks are identical at inference, so this is
a control and not a discovery, and it is the experiment promised by \cref{ch:ltn}.
The second compares the network alone with the network followed by the post-check,
and reports the added latency and the violation rate before and after. It measures
the cost of the symbolic component at inference and bears on RQ4.

% ================================================================
% 7.8 Knowledge base design
% ================================================================
\section{Knowledge base design}
\label{sec:kb-design}

The knowledge base holds the axioms that the knowledge arms use and that every arm
is scored against. It has a format, a way of being written and a version, so that
the axiom set is a factor of the benchmark and not a hidden choice inside a model.

\subsection{Elicitation of rules from Chapter 3}
\label{sec:kb-elicitation}

The source is \cref{sec:domain-knowledge-formalisation}, where each attack family is
paired with an observable condition and a candidate axiom
(\cref{tab:candidate-axioms}); U2R has none. Elicitation turns a candidate into an
axiom that runs on a dataset, in four steps. Where no candidate can be written on a
dataset, the same steps are applied to rules found in its training data, and the
axiom is tagged as such (\cref{sec:kb-format}).

First, the predicates are mapped to columns of the flow table. Most candidates need
a statistic over a window or a source. NSL-KDD ships window statistics as columns
(\cref{sec:time-window-connection-aggregates}), and an axiom can read them. The CIC
datasets hold statistics of single flows, so there a candidate that needs a window or
a source cannot be written in this version. Adding window statistics means defining
the window once for training and streaming (\cref{sec:uc2}). The columns read by each
axiom are listed in \cref{app:axioms}.
Second, the thresholds are fixed, on training data only. For a candidate of
\cref{ch:attacks-ddos} the condition comes from the mechanism. Its threshold is
either fixed before the attacks are looked at, as for the slow DoS axiom, whose cuts
are percentiles of Normal traffic, or calibrated on the labelled training data, as
for the SYN flood axiom of NSL-KDD. A rule found in the data takes both its
conditions and its thresholds from a single-feature search or a shallow tree, checked,
where the data allows, on a part of the training data that the search did not use.
The test part is never read, since choosing axioms on it would be the data snooping
that \citet{arp2022dosdonts} warn against. Third, the statistics of the axiom on
training data are recorded: the support, the share of target-class flows that the
antecedent covers; the confidence, the share of covered flows that belong to the
target class; and the number of Normal flows covered. They document the choice and
are not results. Fourth, the axiom is accepted or refused. It enters the main
knowledge base of a dataset when its training confidence is at least 0.90. The bar is
the confidence of the weakest axiom of the NSL-KDD base, the first to be written, and
it was fixed before the bases of the CIC datasets. We also exclude an axiom wholly
contained in another and keep one condition for each group of nearly identical
columns, such as packet rate and inter-arrival time. An axiom that fails the bar but
is informative goes into a variant.

Two cases show what happens when the knowledge is weak. The slow DoS axiom, S1, a
long flow with a low packet rate, does not hold on CIC-IDS2017: of the 31\,079
flows on which it fires, 30\,839 are Normal. It is kept in a variant, as degraded
knowledge. For R2L and U2R on NSL-KDD, single-feature rules, rules by class and a
tree of depth three gave no rule that reached the bar, so the main base has no axiom
for them (\cref{sec:uc3}). One limit is also declared. On CIC-DDoS2019 the first
axiom says that a UDP flow with no backward packet is DDoS, and the condition on
backward packets acts as a shortcut, since it covers almost all attack flows. It is
still a column of the flow table that every arm sees, and not an artefact of the
capture, so we keep it; E3, which removes the axiom, shows how much it matters.

\subsection{Versioned knowledge base format, so that axiom sets are an experiment factor}
\label{sec:kb-format}

A knowledge base is one YAML file for one dataset and one version.
\Cref{lst:kb} shows one axiom of the knowledge base of CIC-DDoS2019.

\begin{lstlisting}[caption={The first axiom of the knowledge base \texttt{cicddos2019\_v1}.},label={lst:kb},float=t,numbers=none,basicstyle=\ttfamily\scriptsize,breaklines=true,frame=single]
version: cicddos2019_v1
axioms:
  - id: A1
    group: dos
    origin: mechanism
    target: DDoS
    train: {support: 0.902, confidence: 0.999, benign_hit: 305}
    all:
      - {col: Protocol, op: eq, value: udp}
      - {col: Total Backward Packets, op: le, tau: 0.5, s: 0.1}
\end{lstlisting}

Each axiom has an \texttt{id}, unique in the file, and a \texttt{target} class named
by the consequent. Its \texttt{group} is the kind of knowledge: \texttt{dos} for
rules that state a mechanism of the DoS family of \cref{ch:attacks-ddos}, which
includes DDoS, flooding or slow; \texttt{dataset} for rules tied to the features of
one dataset; \texttt{domain} for general network knowledge, not used in this version.
The group lets E3 remove a whole kind of knowledge at once. Its \texttt{origin} says
where the rule comes from: \texttt{mechanism} when its condition states a mechanism
of \cref{ch:attacks-ddos}, whatever the source of the threshold; \texttt{data} when
conditions and thresholds were found by a search on training data; and
\texttt{internship} for axioms taken over from the author's earlier project and not
re-derived from the benchmark data. The field \texttt{train} holds the statistics of
the previous section for documentation, and no model reads it. The field
\texttt{all} lists conditions that must all hold. A condition names a column and a
test: equality to a value, or greater than or at most a threshold $\tau$ with a
width $s$.

The axiom reads as the implication of \cref{ch:ltn}: if every condition holds, the
flow belongs to the target class. An equality is crisp. A comparison is graded, with
degree $\sigma((v-\tau)/s)$ for ``greater than'' and the symmetric form for ``at
most'', or crisp when the run asks for hard thresholds. The degree of the antecedent
is the product of the degrees of its conditions. The width $s$ is fixed and not
learned, as in \cref{ch:ltn}; a run can scale all widths by one factor, so the
widths in the file are starting values.

\Cref{tab:kb-versions} in \cref{app:axioms} lists the knowledge bases of the thesis.
Each variant contains its base and adds axioms. \texttt{\_r2l}, \texttt{\_slow} and
\texttt{\_e9ssh} add one axiom each to \texttt{cicids2017\_v1}. \texttt{\_weak} adds to
\texttt{nslkdd\_v1} one axiom for each term of the two disjunctions of
\cref{sec:axiom-set-internship}, since an axiom of the file is a conjunction; U2
covers no U2R flow of the training part and is kept as the disjunction has it. Of
the ten axioms of the other bases, four state a mechanism and six were found in the
data. Only DoS and Probe are covered on NSL-KDD outside the weak variant, and only
DDoS on CIC-DDoS2019, whose knowledge base has a single target class, so the
comparison of knowledge on rare classes cannot be made there (\cref{sec:uc2}). The
three knowledge bases are not the same set. An experiment that moves a model between
datasets (E9) needs the axioms computable on both, or must declare that the
knowledge base changes. The transfer of E9 uses the two axioms of
\texttt{cicids2017\_v1} written in the common columns, and a variant that adds the UDP
axiom of \texttt{cicddos2019\_v1}.

The knowledge base is also a factor of the benchmark. A run records the version of
the knowledge base, the identifiers of its axioms, any edit applied to it, the weight
of the domain axioms, whether conditions were crisp, the exponent $p$, the scale of
the widths, and the mean degree of each antecedent on the training data, which shows
whether an axiom fires at all. A table of results can then be traced to the axiom
set that produced it. The levels of the factor are of two kinds. A version is a
file, and a file is not edited once it is used: the variants of
\cref{tab:kb-versions} are versions. An edit is derived from a version when the run
starts and recorded with the run, while the file stays as it is. The axioms are
always scored against the file, so the violation rates of edited runs stay
comparable. The experiments use the factor in five ways.
\begin{itemize}
  \item E3 removes one axiom, or a whole group, from a version.
  \item E11 degrades the knowledge by moving every threshold of a version by a
    level, measured in percentile points of the training column, or by swapping the
    target of one axiom. An equality has no threshold and changes only through the
    swap. We use levels and not a fraction of corrupted axioms, because the main
    bases hold two or three axioms and a fraction would take at most four values.
  \item E3 and E11 also use \texttt{nslkdd\_v1\_weak}, with broad axioms of low
    precision for R2L and U2R, to test whether such knowledge helps or harms where
    no precise rule exists.
  \item E13 adds two to eight copies of one axiom, with thresholds moved by small
    steps on either side. The cost of a training step does not depend on what an
    axiom says, so copies measure cost against the number of axioms. The use of
    copies is declared where the result is read.
  \item E16 estimates the effort of writing each version and of changing one rule,
    beside the number of axioms and lines of each version, which are counted.
\end{itemize}

The settings of the logic belong to the run, not to the file: the weight of the
domain axioms in the loss, the exponent $p$, the scale of the widths, and crisp or
graded conditions. Experiment E4 varies the exponent $p$, graded against crisp
conditions and the scale of the widths. The connectives, product and Reichenbach,
and the stabilising constant $\epsilon$ of the aggregator are fixed
(\cref{ch:ltn}). $\epsilon$ is a numerical guard with a small value, and changing
it would alter the numerics more than the semantics.

Every arm is scored against a knowledge base, trained on it or not, so that
axiom-level satisfaction and the domain-rule violation rate (C5) are comparable
between arms (FR7). A flow violates an axiom when the antecedent holds and the
predicted class differs from the target. The violation is defined relative to the
rule and not to the label, so a rule that is wrong on a flow is violated by the
correct answer. This is why the precision of the rule is reported beside the
violation rate.

% ================================================================
% 7.9 Summary
% ================================================================
\section{Summary}
\label{sec:ch7-summary}

This chapter fixed the use cases, the requirements and the design that the rest of
the thesis builds on. UC1, UC2 and UC3 are the tasks on which the arms are scored:
multi-class classification of flows, binary DDoS detection, and the rare classes
R2L and U2R, for which we expect a negative result for the knowledge on NSL-KDD.
UC4 and UC5 ask the same detectors for robustness to unseen attacks and for
alerts an analyst can read. \Cref{tab:functional-requirements,tab:nonfunctional-requirements}
give each requirement an identifier, an experiment from E1 to E17 and a measure on
one of the three axes.

Two pipelines share a model, a feature definition and a knowledge base. In batch
the symbolic component acts at training only, and at inference an LTN-trained
network is a plain network; an optional post-check evaluates the axioms on the raw
flow, to annotate or to enforce. The streaming pipeline is specified, and its
measurement depends on infrastructure that is not yet ours. The knowledge base is a
versioned file, and its versions and recorded edits are the levels of an
experimental factor.

\Cref{ch:benchmarking-methodology} takes this design and fixes how it is
measured: the datasets and their splits, the arms and their tuning, the measures
and the statistical tests.
